% =====================================================================
% The packages, the macros and the conventions of this thesis.
%
% The class loads fontspec, geometry, setspace, polyglossia and bidi
% itself, so this file only adds what the chapters need. bidi is
% sensitive to the order in which packages are loaded, so anything
% added here is loaded after the class.
% =====================================================================

% --------------------------- general --------------------------------
\usepackage{xspace}
\usepackage{xcolor}
\usepackage{soul}
\usepackage{ulem}
\normalem

% ---- review markup (Dana's review of Neta's changes) ----
\definecolor{reviewblue}{HTML}{1565C0}
\definecolor{reviewpink}{HTML}{C62828}
\newcommand{\dremove}[1]{{\color{reviewblue}\sout{#1}}}
\newcommand{\dadd}[1]{{\color{reviewpink}#1}}
\newcommand{\Dana}[1]{\par\noindent{\color{reviewpink}\textbf{Dana:} #1}\par}
\definecolor{revieworange}{HTML}{E65100}
\newcommand{\Danawarn}[1]{\par\noindent{\color{revieworange}\textbf{Dana (open):} #1}\par}
% pinned to the 2024-06-01 release. array 2.7 dropped
% \UseMathForPositioningText, which bidi's right-to-left patch for
% tabular still calls, so with the current release every table in the
% thesis fails. Drop the date once bidi has caught up.
\usepackage{array}[=2024-06-01]
\usepackage{graphicx}
\usepackage[labelfont=bf]{caption}
\usepackage{subcaption}
\usepackage{tikz}
\usepackage{hyperref}
\hypersetup{hidelinks,linktoc=all}

% ---------------------- formulas and equations -----------------------
\usepackage{amsmath}
\usepackage{amssymb}
\usepackage{amsthm}
\usepackage{mathtools}

% equations numbered within chapters: (2.1), (2.2), ...
\numberwithin{equation}{chapter}

\theoremstyle{plain}
\newtheorem{theorem}{Theorem}[chapter]
\newtheorem{lemma}[theorem]{Lemma}
\newtheorem{proposition}[theorem]{Proposition}
\newtheorem{corollary}[theorem]{Corollary}

\theoremstyle{definition}
\newtheorem{definition}[theorem]{Definition}
\newtheorem{example}[theorem]{Example}

\theoremstyle{remark}
\newtheorem{remark}[theorem]{Remark}

% the notation, with the names of the paper, so that formulas carry over
\newcommand{\reals}{\mathbb{R}}
\newcommand{\relu}{\mathrm{ReLU}}
\newcommand{\Npre}{N_{\mathrm{pre}}}
\newcommand{\ep}{\varepsilon}
\newcommand{\mI}{\mathcal{I}}
\newcommand{\conf}[3]{\mathcal{C}(#1,#2,#3)}  % network, input, class
\newcommand{\dlo}{\delta_l}
\newcommand{\dup}{\delta_u}

% The chapters were written with the citation commands of natbib;
% biblatex, which the class uses, spells the same two commands
% \parencite and \textcite.
\let\cite\parencite
\let\citep\parencite
\let\citet\textcite

% Some of the venue names in the bibliography are long and carry no
% break point, so the entries are set with relaxed spacing.
\appto{\bibsetup}{\sloppy}

% the names of the tool and of the baseline
\newcommand{\tool}{BLEND\xspace}
\newcommand{\baseline}{VHAGaR\xspace}

% ---------------------------- figures --------------------------------
\graphicspath{{graphics/}}
\usetikzlibrary{arrows.meta,calc,positioning,patterns,shapes.geometric}
\captionsetup[figure]{position=below,skip=6pt}

% The legend and the column headings of the tau figure are cropped off
% the drawing and set here instead. The drawing is a finished pdf from
% the paper, with everything in it at one size, so scaling cannot
% enlarge its labels alone; set in the document they read at \small
% rather than at the 8.2pt the drawing gives them. The colours are the
% ones the drawing uses for the two series.
% the gutter the row labels of the tau figure sit in, and what is
% left of the text block for the drawing itself
\newlength{\taukeywd}
\newlength{\tauimgwd}
\definecolor{blendplain}{HTML}{1F78B4}
\definecolor{blendxfer}{HTML}{2CA02C}
\newcommand{\serieskey}[2]{%
  \tikz[baseline=-0.55ex]{%
    \draw[#1,line width=0.7pt] (0,0) -- (12pt,0);%
    \node[#2,fill=#1,draw=#1,inner sep=1.3pt] at (6pt,0) {};}}

% ---------------------------- tables ---------------------------------
\usepackage{booktabs}
\usepackage{multirow}
\usepackage{longtable}
\usepackage{siunitx}
\usepackage{adjustbox}   % used by the result tables copied from the paper
\captionsetup[table]{position=above,skip=6pt}

% The evaluation sets large tables with [H], which forces a page break
% whenever one does not fit in what is left of the page. Under
% \flushbottom, which report gives a two-sided document, the glue on
% such a page is stretched until the page is full, which drives the
% heading to the top and its paragraph to the foot with a hole between
% them. \raggedbottom leaves the slack at the bottom instead, where it
% belongs.
\raggedbottom

% [H], which pins a table exactly where it is written instead of
% letting it float. Used for the small tables of the evaluation, which
% otherwise sink to the foot of the page, away from the paragraph that
% introduces them. Only for tables that are short enough to fit in the
% space left on the page; a large one would leave a gap instead.
\usepackage{float}

% \FloatBarrier, used at the end of every subsection of the evaluation
% to keep a table inside the subsection that discusses it. Without it a
% float can drift pages past its text, because floats are emitted in
% the order they are declared and one stuck float holds up the rest.
\usepackage{placeins}

% --------------------------- the algorithm ---------------------------
\usepackage[ruled,linesnumbered,noend]{algorithm2e}
\DontPrintSemicolon
\SetKwFunction{BLEND}{BLEND}
\SetKwFunction{PGD}{PGDAttack}

% ------------------- the filler pages are blank ----------------------
% The class prints the thesis two-sided, so that every chapter opens on
% a right-hand page, and LaTeX inserts a filler page wherever that
% leaves a left-hand page empty. Those pages are meant to be blank, but
% the page style in force still prints a number on them. These are the
% two commands that insert them, with the filler page made empty.
\makeatletter
\renewcommand{\cleardoublepage}{%
  \clearpage
  \if@twoside
    \ifodd\c@page\else
      \hbox{}\thispagestyle{empty}\newpage
      \if@twocolumn\hbox{}\newpage\fi
    \fi
  \fi}
\renewcommand*{\cleartoleftpage}{%
  \clearpage
  \if@twoside
    \ifodd\c@page
      \hbox{}\thispagestyle{empty}\newpage
      \if@twocolumn\hbox{}\newpage\fi
    \fi
  \fi}
\makeatother

% ---------------- the width of the notation list ---------------------
% The style of the class gives the symbols a column 3cm wide, which the
% difference intervals overrun, and its columns add up to more than the
% text block, because the padding between them is not subtracted. This
% is the same style with a wider first column and the arithmetic fixed.
% Its third column lists the pages an entry is used on; the entries
% here are never cited with \gls, so that column is always empty and
% is given no width, which would otherwise leave the list some 50pt
% narrower than the text block.
\renewglossarystyle{iitthesis}{%
    \renewenvironment{theglossary}%
    {\begin{longtable}{@{}p{4.6cm}@{\quad}%
       p{\dimexpr\textwidth-4.6cm-1em-2\tabcolsep\relax}%
       p{0pt}@{}}}%
        {\end{longtable}}%
    \renewcommand*{\glossaryheader}{}%
    \renewcommand*{\glsgroupheading}[1]{}%
    \renewcommand*{\glossentry}[2]{ %
        \textbf{\glstarget{##1}{\glossentryname{##1}}}%
        &
        \formatglossarydescription{##1}%
        & ##2 %
        \tabularnewline%
    }%
    \renewcommand*{\glsgroupskip}{}%
}

% ------------------- one list of figures and tables ------------------
% The class prints a list of figures alone. The thesis has twelve
% result tables as well, so the two lists are printed together, under
% one heading, where the class would have printed the list of figures.
\makeatletter
\renewcommand{\conditionallistoffigures}{%
  \@ifundefined{r@iitthesis:had-figures}{}{%
    \chapter*{List of Figures and Tables}%
    \addcontentsline{toc}{chapter}{List of Figures and Tables}%
    \addcontentsline{cot}{chapter}{רשימת איורים וטבלאות}%
    \@starttoc{lof}%
    \@starttoc{lot}%
  }%
}
\makeatother
